% Auto-generated from meeting_2568-10-11_admission-committee.md (กบ.วช. format v2)
\documentclass{meetingminutes}
\meetingcommittee{คณะกรรมการบริหารหลักสูตร วิศวกรรมคอมพิวเตอร์และปัญญาประดิษฐ์ (บัณฑิตศึกษา)}
\meetingnumber{๑}{๒๕๖๘}
\meetingdate{วันที่ ๑๑ ตุลาคม พ.ศ. ๒๕๖๘}
\meetingplace{ห้องประชุมคณะเทคโนโลยีอุตสาหกรรม}
\meetingstart{๐๙.๐๐ น.}
\meetingend{๑๐.๕๐ น.}

\begin{document}
\maketitle

\psection{ผู้มาประชุม}
\begin{participants}
  \pmember{1}{ผู้ช่วยศาสตราจารย์ ดร.พิศิษฐ์ นาคใจ}{อาจารย์ประจำหลักสูตร}{ประธาน}
  \pmember{2}{ผู้ช่วยศาสตราจารย์ ดร.กาญจนา ดาวเด่น}{อาจารย์ประจำหลักสูตร (เลขานุการ)}{กรรมการ/เลขานุการ}
  \pmember{3}{ผู้ช่วยศาสตราจารย์ ดร.พัชรี มณีรัตน์}{อาจารย์ประจำหลักสูตร}{กรรมการ}
\end{participants}

\psection{ผู้ไม่มาประชุม}
\noindent ไม่มี\par

\startmeeting
\openingchair{ผู้ช่วยศาสตราจารย์ ดร.พิศิษฐ์ นาคใจ}{กล่าวเปิดการประชุมและดำเนินการประชุมตามระเบียบวาระ ดังนี้}

\agenda{๑}{เรื่องแจ้งให้ที่ประชุมทราบ}
ประธานแจ้งที่ประชุมว่า หลักสูตร วศ.ม. วิศวกรรมคอมพิวเตอร์และปัญญาประดิษฐ์ เริ่มเปิดรับสมัครนักศึกษารุ่นแรกอย่างเป็นทางการแล้ว โดยมีผู้สนใจติดต่อสอบถามเข้ามาอย่างต่อเนื่องผ่านช่องทาง Facebook Page และเว็บไซต์ของมหาวิทยาลัย


\agenda{๒}{กำหนดการและเกณฑ์การรับสมัคร}
ที่ประชุมยืนยันเกณฑ์การรับสมัครดังนี้:

เกณฑ์คุณสมบัติผู้สมัคร: 1. สำเร็จการศึกษาระดับปริญญาตรีหรือเทียบเท่าในสาขาวิศวกรรมคอมพิวเตอร์ วิทยาการคอมพิวเตอร์ เทคโนโลยีสารสนเทศ หรือสาขาที่เกี่ยวข้อง 2. มีคะแนนเฉลี่ยสะสม (GPA) ไม่ต่ำกว่า 2.50 หรือมีประสบการณ์ทำงานที่เกี่ยวข้อง 3. ผ่านการสอบสัมภาษณ์จากคณะกรรมการ

กำหนดการ: - รับสมัคร: 1 ตุลาคม – 15 พฤศจิกายน 2568 - สอบสัมภาษณ์: กลางพฤศจิกายน 2568 - ประกาศผล: ปลายพฤศจิกายน 2568 - รายงานตัว: ต้นธันวาคม 2568 - เปิดภาคเรียน: พฤศจิกายน 2568

\resolution{เห็นชอบเกณฑ์และกำหนดการรับสมัครตามที่เสนอ}

\agenda{๓}{กำหนดขั้นตอนการสอบสัมภาษณ์}
ที่ประชุมกำหนดขั้นตอนการสอบสัมภาษณ์: 1. คณะกรรมการสัมภาษณ์: ประกอบด้วยอาจารย์ประจำหลักสูตรอย่างน้อย 3 คน 2. ประเด็นการสัมภาษณ์: ความรู้พื้นฐาน, แรงจูงใจในการศึกษาต่อ, แผนการวิจัย/สารนิพนธ์ 3. เกณฑ์ผ่าน: ผลการสัมภาษณ์ 60\% ขึ้นไป

\resolution{เห็นชอบขั้นตอนการสอบสัมภาษณ์ และมอบหมายให้ ผศ.ดร.พิศิษฐ์ เป็นประธานคณะกรรมการสัมภาษณ์}

\adjournmeeting

\begin{sigblock}
\typist{ผู้ช่วยศาสตราจารย์ ดร.กาญจนา ดาวเด่น}
\reviewer{ผู้ช่วยศาสตราจารย์ ดร.พิศิษฐ์ นาคใจ}{กรรมการ}
\chairsig{ผู้ช่วยศาสตราจารย์ ดร.พิศิษฐ์ นาคใจ}{ประธานการประชุม}
\end{sigblock}

\end{document}
